\documentclass[11pt]{article}
\usepackage[margin=1in]{geometry}
\usepackage{xcolor}
\usepackage[breakable,listings]{tcolorbox}
\tcbuselibrary{listings}
\usepackage{hyperref}
\usepackage{enumitem}
\usepackage{listings}
\usepackage{verbatim}
\usepackage{amsmath}
\usepackage{graphicx}
\usepackage{booktabs}
\usepackage{float}

% Define colors
\definecolor{osaorange}{RGB}{220,115,38}
\definecolor{codebackground}{RGB}{248,248,248}
\definecolor{answergreen}{RGB}{0,128,0}

% Configure hyperref
\hypersetup{
    colorlinks=true,
    linkcolor=blue,
    urlcolor=blue,
    citecolor=blue
}

% Configure listings for code display
\lstset{
  basicstyle=\small\ttfamily,
  breaklines=true,
  breakatwhitespace=true,
  postbreak=\mbox{\textcolor{gray}{$\hookrightarrow$}\space},
  columns=flexible,
  keepspaces=true,
  showstringspaces=false
}

% Code environment using tcolorbox with listings
\newtcblisting{rcode}{
  colback=white,
  colframe=gray!50,
  boxrule=0.5pt,
  arc=0pt,
  left=3pt,
  right=3pt,
  top=3pt,
  bottom=3pt,
  width=\textwidth,
  breakable,
  listing only,
  listing options={
    basicstyle=\small\ttfamily,
    breaklines=true,
    breakatwhitespace=true,
    postbreak=\mbox{\textcolor{gray}{$\hookrightarrow$}\space},
    columns=flexible,
    keepspaces=true,
    showstringspaces=false
  }
}

\title{}
\author{}
\date{}

\begin{document}

% Header box
\begin{tcolorbox}[colback=osaorange,colframe=black,boxrule=2pt,arc=0pt,left=3pt,right=3pt,top=6pt,bottom=6pt,boxsep=2pt]
\centering
\textcolor{white}{\textbf{\Large H524 Introduction to Biostatistics}}\\
\textcolor{white}{\textbf{\Large Week 7 Lab: Nonparametric Methods - ANSWER KEY}}\\
\textcolor{white}{\textbf{\Large Fall 2025}}
\end{tcolorbox}

\vspace{8pt}

\section{Exercise 1: Shapiro-Wilk Test}

\begin{rcode}
# Generate example data
set.seed(524)
normal_data <- rnorm(30, mean=100, sd=15)
skewed_data <- rexp(30, rate=0.02)

# Check normality
shapiro.test(normal_data)
shapiro.test(skewed_data)
\end{rcode}

\textbf{Answers:}
\begin{enumerate}
\item \textcolor{answergreen}{P-value for \texttt{normal\_data}: 0.2371}
\item \textcolor{answergreen}{P-value for \texttt{skewed\_data}: 0.0025}
\item \textcolor{answergreen}{For \texttt{normal\_data}, use \textbf{parametric} tests (p = 0.2371 $\geq$ 0.05, normality assumption OK)}
\item \textcolor{answergreen}{For \texttt{skewed\_data}, use \textbf{nonparametric} tests (p = 0.0025 < 0.05, data not normal)}
\end{enumerate}

\section{Exercise 2: Sign Test for One Sample}

\textbf{Question:} Is median recovery time different from 7 days?

\begin{rcode}
# Recovery times (days)
recovery <- c(8, 5, 7, 10, 4, 9, 3, 12, 6, 11, 7)

# Hypothesized median
med0 <- 7

# Calculate differences
diff <- recovery - med0

# Count signs (ignoring zeros)
n_positive <- sum(diff > 0)
n_total <- sum(diff != 0)

cat("Positive differences:", n_positive, "\n")
cat("Total non-zero:", n_total, "\n")

# Sign test
binom.test(n_positive, n_total, p=0.5)
\end{rcode}

\textbf{Answers:}
\begin{enumerate}
\item \textcolor{answergreen}{Positive differences: 5}
\item \textcolor{answergreen}{Total non-zero differences: 9}
\item \textcolor{answergreen}{P-value: 1.0}
\item \textcolor{answergreen}{Is median significantly different from 7 days? \textbf{NO} (p = 1.0 $\geq$ 0.05)}
\end{enumerate}

\section{Exercise 3: Sign Test for Paired Data}

\textbf{Question:} Does treatment reduce pain scores?

\begin{rcode}
# Pain scores (1-10 scale)
before <- c(7, 8, 6, 9, 5, 8, 7, 6)
after <- c(5, 6, 4, 7, 4, 6, 5, 5)

# Calculate differences
diff <- before - after

# Count signs
n_improved <- sum(diff > 0)
n_total <- sum(diff != 0)

cat("Number improved:", n_improved, "\n")
cat("Number total:", n_total, "\n")

# Sign test
binom.test(n_improved, n_total, p=0.5)

# Summary
mean(before)
mean(after)
mean(diff)
\end{rcode}

\textbf{Answers:}
\begin{enumerate}
\item \textcolor{answergreen}{Number improved: 8}
\item \textcolor{answergreen}{P-value: 0.0078}
\item \textcolor{answergreen}{Does treatment reduce pain? \textbf{YES} (p = 0.0078 < 0.05)}
\item \textcolor{answergreen}{Average reduction in pain score: 1.75}
\end{enumerate}

\section{Exercise 4: Blood Pressure Study}

\textbf{Research Question:} Does medication reduce blood pressure?

\begin{rcode}
# Blood pressure before and after medication
bp_before <- c(145, 138, 150, 142, 156, 148, 140, 152)
bp_after <- c(138, 135, 142, 140, 148, 145, 138, 146)

# Check normality of differences
diff <- bp_before - bp_after
shapiro.test(diff)

# Wilcoxon signed-rank test
wilcox.test(bp_before, bp_after, paired=TRUE)

# Summary statistics
mean(bp_before)
mean(bp_after)
mean(diff)
\end{rcode}

\textbf{Answers:}
\begin{enumerate}
\item \textcolor{answergreen}{Shapiro-Wilk p-value for differences: 0.0711}

\textcolor{answergreen}{Normality assumption OK? YES (p = 0.0711 $\geq$ 0.05, though borderline)}

\item \textcolor{answergreen}{Wilcoxon test p-value: 0.0139}

\item \textcolor{answergreen}{Significant at $\alpha = 0.05$? YES (p = 0.0139 < 0.05)}

\item \textcolor{answergreen}{Does medication reduce BP? YES ($p < 0.05$, significant reduction)}

\item \textcolor{answergreen}{Mean reduction in BP: 4.88 mmHg}
\end{enumerate}

\section{Exercise 5: Physical Therapy Study}

\textbf{Question:} Does physical therapy reduce pain scores?

\begin{rcode}
# Pain scores (0-10 scale) before and after 8 weeks PT
pain_before <- c(7, 8, 6, 9, 7, 8, 6, 7, 8, 9)
pain_after <- c(5, 6, 5, 7, 6, 7, 4, 6, 7, 8)

# Check normality
diff <- pain_before - pain_after
shapiro.test(diff)

# Wilcoxon test
wilcox.test(pain_before, pain_after, paired=TRUE)

# Medians
median(pain_before)
median(pain_after)
median(diff)
\end{rcode}

\textbf{Answers:}
\begin{enumerate}
\item \textcolor{answergreen}{Are differences normally distributed? \textbf{NO} (p = 0.0002 < 0.05)}
\item \textcolor{answergreen}{Wilcoxon p-value: 0.0046}
\item \textcolor{answergreen}{Does PT significantly reduce pain? \textbf{YES} ($p < 0.05$)}
\item \textcolor{answergreen}{Median reduction in pain score: 1}
\end{enumerate}

\section{Exercise 6: Treatment Comparison}

\textbf{Research Question:} Do two treatments differ in recovery time?

\begin{rcode}
# Recovery times (days)
treatment_A <- c(23, 31, 25, 28, 30, 27, 29)
treatment_B <- c(18, 22, 20, 24, 19, 21, 23)

# Check normality
shapiro.test(treatment_A)
shapiro.test(treatment_B)

# Wilcoxon rank-sum test
wilcox.test(treatment_A, treatment_B)

# Summary statistics
median(treatment_A)
median(treatment_B)
\end{rcode}

\textbf{Answers:}
\begin{enumerate}
\item \textcolor{answergreen}{Treatment A normal? YES (p = 0.8525 $\geq$ 0.05)}

\textcolor{answergreen}{Treatment B normal? YES (p = 0.9493 $\geq$ 0.05)}

\textcolor{answergreen}{Note: Both are normal, but we can still use Wilcoxon as a robust alternative}

\item \textcolor{answergreen}{Wilcoxon rank-sum p-value: 0.004}

\item \textcolor{answergreen}{Significant at $\alpha = 0.05$? YES (p = 0.004 < 0.05)}

\item \textcolor{answergreen}{Which treatment has shorter recovery? \textbf{Treatment B} (median 21 days vs 28 days)}

\item \textcolor{answergreen}{Median recovery times:}
\begin{itemize}
\item \textcolor{answergreen}{Treatment A: 28 days}
\item \textcolor{answergreen}{Treatment B: 21 days}
\end{itemize}
\end{enumerate}

\section{Exercise 7: Drug Effectiveness Study}

\textbf{Question:} Do two drugs differ in effectiveness?

\begin{rcode}
# Symptom relief scores (higher = better)
drug_X <- c(45, 52, 48, 55, 50, 47, 53)
drug_Y <- c(38, 42, 40, 45, 39, 43, 41)

# Wilcoxon test
wilcox.test(drug_X, drug_Y)

# Summary
mean(drug_X)
mean(drug_Y)
median(drug_X)
median(drug_Y)
\end{rcode}

\textbf{Answers:}
\begin{enumerate}
\item \textcolor{answergreen}{Wilcoxon p-value: 0.0026}
\item \textcolor{answergreen}{Which drug is more effective? \textbf{Drug X} (higher scores)}
\item \textcolor{answergreen}{Difference in average scores: 50 - 41.14 = 8.86 points}
\end{enumerate}

\section{Exercise 8: Pain Relief Comparison}

\textbf{Research Question:} Do three pain medications differ in effectiveness?

\begin{rcode}
# Pain relief scores (0-10 scale, higher = better)
drug_A <- c(3, 4, 5, 4, 3, 5, 4)
drug_B <- c(5, 6, 7, 6, 5, 7, 6)
drug_C <- c(7, 8, 9, 8, 7, 9, 8)

# Combine data
pain_scores <- c(drug_A, drug_B, drug_C)
drug <- factor(rep(c("A", "B", "C"), each=7))

# Kruskal-Wallis test
kruskal.test(pain_scores ~ drug)

# Post-hoc pairwise comparisons
pairwise.wilcox.test(pain_scores, drug, p.adjust.method="bonferroni")

# Summary statistics
tapply(pain_scores, drug, median)
\end{rcode}

\textbf{Answers:}
\begin{enumerate}
\item \textcolor{answergreen}{Kruskal-Wallis p-value: 0.0002}

\item \textcolor{answergreen}{Significant overall at $\alpha = 0.05$? YES (p = 0.0002 < 0.05)}

\item \textcolor{answergreen}{Post-hoc tests (Bonferroni-adjusted):}
\begin{itemize}
\item \textcolor{answergreen}{A vs B: p = 0.0121 (SIGNIFICANT)}
\item \textcolor{answergreen}{A vs C: p = 0.0057 (SIGNIFICANT)}
\item \textcolor{answergreen}{B vs C: p = 0.0121 (SIGNIFICANT)}
\item \textcolor{answergreen}{All three drugs differ significantly from each other}
\end{itemize}

\item \textcolor{answergreen}{Most effective drug: \textbf{Drug C} (median = 8)}

\item \textcolor{answergreen}{Least effective drug: \textbf{Drug A} (median = 4)}
\end{enumerate}

\section{Exercise 9: Comparing Parametric vs Nonparametric}

\textbf{Question:} Sometimes data are borderline normal. Let's compare both approaches.

\begin{rcode}
# Blood pressure data (from Exercise 4)
bp_before <- c(145, 138, 150, 142, 156, 148, 140, 152)
bp_after <- c(138, 135, 142, 140, 148, 145, 138, 146)

# Check normality
diff <- bp_before - bp_after
shapiro.test(diff)

# Parametric: Paired t-test
t_result <- t.test(bp_before, bp_after, paired=TRUE)
cat("Paired t-test p-value:", t_result$p.value, "\n")

# Nonparametric: Wilcoxon
w_result <- wilcox.test(bp_before, bp_after, paired=TRUE)
cat("Wilcoxon p-value:", w_result$p.value, "\n")

# Compare
cat("\nBoth tests significant? Compare p-values to 0.05\n")
\end{rcode}

\textbf{Answers:}
\begin{enumerate}
\item \textcolor{answergreen}{Is normality assumption met ($p \geq 0.05$)? \textbf{YES} (borderline, p = 0.0711)}
\item \textcolor{answergreen}{Which test gives smaller p-value? \textbf{t-test} (0.0012 vs 0.0139)}
\item \textcolor{answergreen}{Do both reach same conclusion? \textbf{YES} (both significant, $p < 0.05$)}
\item \textcolor{answergreen}{When borderline normal, which is safer? \textbf{Nonparametric} (more robust)}
\end{enumerate}

\section{Practice Problem 1: Hospital Stay Duration}

\begin{rcode}
# Hospital stay (days)
procedure_1 <- c(3, 4, 5, 3, 6, 4, 7, 5, 4, 12)
procedure_2 <- c(5, 6, 8, 7, 9, 6, 8, 10, 7, 15)

# Check normality
shapiro.test(procedure_1)
shapiro.test(procedure_2)

# Wilcoxon rank-sum test
wilcox.test(procedure_1, procedure_2)

# Summary statistics
median(procedure_1)
median(procedure_2)
\end{rcode}

\textbf{Answers:}
\begin{enumerate}
\item \textcolor{answergreen}{Procedure 1 normal? \textbf{NO} (p = 0.0083 < 0.05)}

\textcolor{answergreen}{Procedure 2 normal? \textbf{YES} (p = 0.0538 $\geq$ 0.05, borderline)}

\textcolor{answergreen}{Since Procedure 1 is not normal, use nonparametric test}

\item \textcolor{answergreen}{Wilcoxon rank-sum p-value: 0.0108}

\item \textcolor{answergreen}{Significant difference? \textbf{YES} (p = 0.0108 < 0.05)}

\item \textcolor{answergreen}{Shorter stays: \textbf{Procedure 1} (median 4.5 days vs 7.5 days)}
\begin{itemize}
\item \textcolor{answergreen}{Median Procedure 1: 4.5 days}
\item \textcolor{answergreen}{Median Procedure 2: 7.5 days}
\end{itemize}

\item \textcolor{answergreen}{Why is Wilcoxon appropriate? Data not normal + outliers present (12, 15 days)}
\end{enumerate}

\section{Practice Problem 2: Diet Program Comparison}

\begin{rcode}
# Weight loss (kg) after 3 months
diet_A <- c(2.1, 3.2, 2.5, 3.0, 2.8, 2.6)
diet_B <- c(3.5, 4.2, 3.8, 4.0, 3.9, 3.7)
diet_C <- c(4.5, 5.1, 4.8, 5.3, 4.9, 5.0)
diet_D <- c(2.8, 3.0, 2.9, 3.1, 2.7, 2.9)

# Combine data
weight_loss <- c(diet_A, diet_B, diet_C, diet_D)
diet <- factor(rep(c("A", "B", "C", "D"), each=6))

# Kruskal-Wallis test
kruskal.test(weight_loss ~ diet)

# Post-hoc tests
pairwise.wilcox.test(weight_loss, diet, p.adjust.method="bonferroni")

# Summary statistics
tapply(weight_loss, diet, median)
\end{rcode}

\textbf{Answers:}
\begin{enumerate}
\item \textcolor{answergreen}{Kruskal-Wallis p-value: 0.000195}

\item \textcolor{answergreen}{Significant overall? \textbf{YES} (p = 0.000195 < 0.05)}

\item \textcolor{answergreen}{Post-hoc tests (Bonferroni-adjusted) - Significant pairs:}
\begin{itemize}
\item \textcolor{answergreen}{A vs B: p = 0.013 (SIGNIFICANT)}
\item \textcolor{answergreen}{A vs C: p = 0.013 (SIGNIFICANT)}
\item \textcolor{answergreen}{A vs D: p = 1.000 (NOT significant)}
\item \textcolor{answergreen}{B vs C: p = 0.013 (SIGNIFICANT)}
\item \textcolor{answergreen}{B vs D: p = 0.030 (SIGNIFICANT)}
\item \textcolor{answergreen}{C vs D: p = 0.030 (SIGNIFICANT)}
\end{itemize}

\item \textcolor{answergreen}{Most effective diet: \textbf{Diet C} (median = 4.95 kg)}

\item \textcolor{answergreen}{Median weight loss by diet:}
\begin{itemize}
\item \textcolor{answergreen}{Diet A: 2.7 kg}
\item \textcolor{answergreen}{Diet B: 3.85 kg}
\item \textcolor{answergreen}{Diet C: 4.95 kg}
\item \textcolor{answergreen}{Diet D: 2.9 kg}
\end{itemize}
\end{enumerate}

\section{Practice Problem 3: Satisfaction Ratings}

\textbf{Question:} Compare customer satisfaction ratings across three stores.

\begin{rcode}
# Satisfaction ratings (1-5 scale, ordinal)
store_A <- c(3, 4, 3, 5, 4, 3, 4, 3, 4, 3)
store_B <- c(4, 5, 4, 5, 5, 4, 5, 4, 5, 4)
store_C <- c(2, 3, 2, 4, 3, 2, 3, 2, 3, 3)

# Combine
satisfaction <- c(store_A, store_B, store_C)
store <- factor(rep(c("A", "B", "C"), each=10))

# Kruskal-Wallis (appropriate for ordinal data!)
kruskal.test(satisfaction ~ store)

# Post-hoc
pairwise.wilcox.test(satisfaction, store, p.adjust.method="bonferroni")

# Medians
tapply(satisfaction, store, median)
\end{rcode}

\textbf{Answers:}
\begin{enumerate}
\item \textcolor{answergreen}{Why is nonparametric appropriate? Data is \textbf{ORDINAL} (1-5 rating scale)}
\item \textcolor{answergreen}{Kruskal-Wallis p-value: 0.0002}
\item \textcolor{answergreen}{Do stores differ in satisfaction? \textbf{YES} ($p < 0.05$)}
\item \textcolor{answergreen}{Highest satisfaction: \textbf{Store B} (median = 4.5)}
\item \textcolor{answergreen}{Significant pairs:}
\begin{itemize}
\item \textcolor{answergreen}{A vs B: p = 0.028 (SIGNIFICANT)}
\item \textcolor{answergreen}{A vs C: p = 0.042 (SIGNIFICANT)}
\item \textcolor{answergreen}{B vs C: p = 0.0007 (SIGNIFICANT)}
\end{itemize}
\end{enumerate}

\section{Key Interpretations}

\subsection{Normality Testing}
\begin{itemize}
\item If Shapiro-Wilk $p \geq 0.05$: Data consistent with normal distribution, can use parametric tests
\item If Shapiro-Wilk $p < 0.05$: Data significantly non-normal, use nonparametric tests
\item Borderline cases ($p \approx 0.05$): Use nonparametric for safety
\end{itemize}

\subsection{Wilcoxon Tests}
\begin{itemize}
\item \textbf{Paired:} Use when comparing before/after measurements on same subjects
\item \textbf{Rank-sum:} Use when comparing two independent groups
\item Interpret p-values same as t-tests: $p < 0.05$ = significant difference
\end{itemize}

\subsection{Kruskal-Wallis Test}
\begin{itemize}
\item Tests if at least one group differs from others
\item If significant ($p < 0.05$), follow up with post-hoc pairwise comparisons
\item Bonferroni correction adjusts for multiple comparisons (more conservative)
\end{itemize}

\subsection{Common Student Errors to Watch For}
\begin{enumerate}
\item \textbf{Forgetting to check normality first}
\item \textbf{Using \texttt{paired=TRUE} for independent groups} (should only be for paired data)
\item \textbf{Interpreting medians as means} (nonparametric tests focus on medians)
\item \textbf{Not following up significant Kruskal-Wallis with post-hoc tests}
\item \textbf{Comparing p-values to 0.01 instead of 0.05} (unless specified)
\end{enumerate}

\end{document}
